\introduction

El panorama contemporáneo de las industrias creativas se caracteriza por una alta dependencia de activos digitales. Desde el diseño gráfico y la producción audiovisual hasta el desarrollo de videojuegos, la capacidad de crear, almacenar, localizar y reutilizar recursos digitales de forma eficiente constituye un factor crítico de productividad y competitividad. Sin embargo, la práctica profesional evidencia una gestión fragmentada de estos recursos, con archivos dispersos en múltiples soportes y servicios no integrados. Esta situación provoca pérdida de tiempo, duplicidad de versiones, dificultades de colaboración y disminución del rendimiento de los equipos creativos.

En este contexto, la investigación se orienta al desarrollo de Assets Studio como solución tecnológica para la gestión inteligente de activos digitales. El objeto de estudio se centra en los procesos de organización, recuperación y aprovechamiento de recursos multimedia en entornos con restricciones de conectividad e infraestructura. El campo de acción se delimita al diseño e implementación de una aplicación web con funcionalidades de clasificación, búsqueda semántica, generación asistida por inteligencia artificial e interacción colaborativa.

El diseño metodológico se estructura en dos componentes complementarios. En primer lugar, se aplica Design Thinking para comprender las necesidades de los usuarios, definir el problema y construir prototipos progresivos validados mediante pruebas de usabilidad \citep{Brown2008,Nielsen1994}. En segundo lugar, se adopta Scrum para planificar e implementar iteraciones de desarrollo con incremento funcional verificable en cada ciclo \citep{Schwaber2020}. La combinación de ambos enfoques permite articular una ruta de trabajo centrada en el usuario y, al mismo tiempo, técnicamente controlada.

La investigación emplea métodos teóricos y empíricos. Entre ellos se incluyen análisis documental, entrevistas semiestructuradas, modelado de requisitos, observación de flujos de trabajo y evaluación de resultados del prototipo. Los datos obtenidos en el diagnóstico inicial evidencian que una parte significativa del tiempo laboral se dedica a tareas logísticas de búsqueda y organización de archivos, lo cual justifica la pertinencia del estudio.

La estructura del documento se organiza en tres capítulos. El Capítulo 1 presenta los fundamentos y referentes teórico-metodológicos sobre gestión de activos digitales, experiencia de usuario y metodologías de desarrollo. El Capítulo 2 expone el diseño e implementación de la solución propuesta, así como los artefactos técnicos construidos. El Capítulo 3 desarrolla la validación de la solución, el análisis de resultados y la factibilidad de su aplicación en el contexto estudiado.
 
% \comment{Este es un comentario}.
% \lipsum[1-2]
